% paper/section_3_design.tex
% ─────────────────────────────────────────────────────────────────────────────
% Section 3: Kavach Architecture
% Slot into skeleton.tex where \section{Kavach Architecture} appears.
% Numbers marked [TBD] are filled in from benchmark workstream D output.
% ─────────────────────────────────────────────────────────────────────────────

\section{Kavach Architecture}
\label{sec:design}
\label{sec:architecture}

\subsection{Overview}

Kavach is a runtime guardrail that interposes between an LLM agent orchestrator
and its tool-execution layer. Every tool call proposed by the agent is
intercepted, evaluated by a Parliament of domain-bound Ministers, and sentenced
by a deterministic Speaker before any side-effecting action is permitted to
execute. The end-to-end request path is: user intent is seeded; each proposed
tool call is embedded once; COMPASS scores intent drift; the semantic router
activates the relevant ministers, which retrieve against their corpora in
parallel; and the Speaker aggregates their verdicts under the trajectory
ceiling before returning BLOCK / ESCALATE / ALLOW.

The system exposes a single HTTP endpoint \texttt{POST /hook/parliament}
that the OpenClaw plugin (\S\ref{sec:plugin}) calls synchronously in the
\texttt{pre\_tool\_call} hook. The endpoint returns a JSON verdict
\texttt{\{decision, confidence, ministers\_fired, compass\_score,
drift\_delta\}} within the configured latency budget.

\subsection{Core Components}

\paragraph{Ledger.}
An append-only SQLite journal (\texttt{kavach\_parliament.db}) records every
evaluation request and its verdict. Each row captures:
\texttt{session\_id}, \texttt{turn}, \texttt{tool\_name}, \texttt{action\_embedding}
(768-d float blob), \texttt{verdict}, and \texttt{timestamp\_utc}.
The Ledger is the single source of truth for session state; no in-memory
session object exists, making the server horizontally stateless between
requests. Session risk score is recomputed from the Ledger at each call via
a lightweight SQL aggregate, adding $<$2\,ms overhead.

\paragraph{Semantic Router.}
On each evaluation request the proposed action text is embedded with
\texttt{BAAI/bge-base-en-v1.5} (768-d, cosine similarity). The resulting
vector is compared against five pre-computed domain centroids stored in
\texttt{kavach\_router\_config.json}. Ministers whose centroid cosine similarity
exceeds $\theta_{\text{route}} = 0.40$ are activated; COMPASS always runs
regardless of routing similarity.

The router replaces the stage-bound dispatcher used in early prototypes
(OpenClaw stages Init$\to$Input$\to$Inference$\to$Decision$\to$Execution).
Ministers are \emph{domain-bound}, not stage-bound: EXECUTOR fires on any
action resembling code execution regardless of which pipeline stage proposed it.
This is the key architectural departure from OpenClaw's native trust model
and from contemporaneous work such as NemoClaw~\cite{nvidia2025nemoclaw}, which
enforces a single boundary at the pipeline entry point.

\paragraph{Ministers.}
Each activated Minister performs a semantic scan: the action embedding is
compared via cosine distance against every pattern in the Minister's
ChromaDB collection. A hit is declared when
$\cos(\mathbf{a}, \mathbf{p}) \geq \theta_{\min}$ where
$\theta_{\min} = 0.45$ (Config\,B; see \S\ref{sec:eval}).
Each corpus pattern carries L1/L2/L3 descriptors (intent $\to$ mechanism
$\to$ surface) that are propagated into the verdict log for provenance
(described in \S\ref{sec:provenance}).

The four domain-bound Ministers and their attack-surface coverage:

\begin{itemize}
  \item \textbf{EXECUTOR} — code execution, process injection, persistence
        mechanisms, container escape (100 patterns; MITRE ATT\&CK T1059,
        T1055, T1053, T1611 and sub-techniques).
  \item \textbf{VAULT} — credential access, secret file reads, metadata service
        queries, platform keychain enumeration (100 patterns; T1552, T1528,
        T1555 and sub-techniques).
  \item \textbf{CHANNEL} — data exfiltration via any outbound channel: HTTP/S,
        DNS, SMTP, cloud storage, git push, covert channels (100 patterns;
        T1048, T1567, T1102 and sub-techniques).
  \item \textbf{NAVIGATOR} — session trajectory anomalies: post-task
        persistence, scope widening, recon sequences, anti-forensic actions
        (100 patterns; OWASP ASI01, T1083, T1082, T1070).
\end{itemize}

\paragraph{COMPASS.}
The Calibrated Oracle for Mission-Pattern Alignment and Session Semantics
evaluates \emph{intent alignment} rather than attack-pattern similarity. It
maintains a \emph{session centroid} — an exponential moving average of action
embeddings weighted by turn recency — and measures cosine drift between the
centroid and the seeded user-intent vector (extracted from the first
non-system turn of the session).

\begin{equation}
  \delta_{\text{drift}}(t) =
    1 - \cos\!\left(\bar{\mathbf{c}}_t,\, \mathbf{u}\right)
  \label{eq:drift}
\end{equation}

where $\bar{\mathbf{c}}_t$ is the turn-$t$ session centroid and $\mathbf{u}$ is
the user-intent embedding. When $\delta_{\text{drift}} > \theta_{\text{drift}}$
the call is escalated regardless of individual minister verdicts.

The COMPASS calibration set (\texttt{kavach\_corpus\_v1.json}, key
\texttt{COMPASS}) contains 100 labelled pairs — 50 \textsc{aligned} and 50
\textsc{hijacked} — covering 25 distinct user-intent scenarios each paired with
both a legitimate and an attack-injected proposed action. Youden's $J$
statistic is used to select $\theta_{\text{drift}}$ via
\texttt{compass\_calibrator.py}.

\paragraph{Speaker.}
The Speaker is a deterministic finite function — no learned parameters — that
maps the set of Minister verdicts and COMPASS output to a single sentence:

\begin{equation}
V = \text{Speaker}(\{v_m\}_{m \in M_{\text{active}}},\, \delta_{\text{drift}},\, r_{\text{session}})
\end{equation}

The decision table (applied in priority order):
\begin{enumerate}
  \item Session risk $r_{\text{session}} \geq r_{\text{max}}$ $\Rightarrow$
        \textsc{Block} (session-level trajectory ceiling; fires independently
        of any single Minister's verdict --- this is what makes Kavach a
        session monitor rather than a per-call guard).
  \item Any Minister returns \textsc{Block} $\Rightarrow$ \textsc{Block}.
  \item $\delta_{\text{drift}} > \theta_{\text{drift}}$ $\Rightarrow$ \textsc{Escalate}.
  \item Any Minister returns \textsc{Grey} (\textsc{Escalate}) $\Rightarrow$ \textsc{Escalate}.
  \item All Ministers return \textsc{Allow} $\Rightarrow$ \textsc{Allow}.
\end{enumerate}
The aggregation is a pure veto: a single \textsc{Block} suffices, with no
confidence-weighted averaging. \S\ref{sec:frontier} studies this choice
empirically against Bayesian and hybrid alternatives under adaptive attack.

\paragraph{Bayesian alternative: priors are estimated, not fixed.}
The Bayesian aggregator we compare against in \S\ref{sec:frontier} does not use
hardcoded conditional probabilities. Its prior and per-Minister reliabilities are
maintained as Beta--Bernoulli posteriors over ground-truth-labelled outcomes.
The block prior is the shrinkage estimate
\begin{equation}
\hat{\pi}_{\text{block}} = \alpha\,\frac{b}{N} + (1-\alpha)\,\pi_0,
\qquad \alpha = \min\!\left(\tfrac{N}{100},\,1\right),
\end{equation}
where $b$ is the number of \textsc{Block} ground-truth labels among $N$ labelled
ledger entries and $\pi_0$ a weak default. Each Minister's reliability
$r_{m,v}$ for vote type $v$ is the running posterior mean of a
Beta$(\cdot)$ updated as
\begin{equation}
r_{m,v} \leftarrow \frac{r_{m,v}\,n_{m,v} + \mathbb{1}[\text{vote correct}]}{n_{m,v}+1},
\qquad n_{m,v} \leftarrow n_{m,v}+1 .
\end{equation}
Crucially, these updates are applied \emph{offline} against verified ground
truth and never at inference time. This is a deliberate security decision:
allowing the runtime to fabricate its own ``ground truth'' from underdetermined
self-decisions would create a feedback channel an adversary could poison to
drift the priors. The deployed Speaker is the deterministic veto above; the
Bayesian path is an analyzed alternative, not the enforcement path.

\subsection{Taxonomy-grounded provenance}
\label{sec:provenance}
Every verdict carries a provenance chain that grounds it in a structured
cyber-taxonomy: the matched corpus pattern resolves to a technique identifier
(MITRE ATLAS for AI/agent-native techniques, e.g.\ AML.T0051 indirect prompt
injection and AML.T0086 exfiltration via agent tool invocation; MITRE ATT\&CK
for host-level techniques), which maps to a tactic and an ordered kill-chain
stage. Resolution is corpus-tagged: every one of the 401 patterns declares a
specific source technique (MITRE ATT\&CK, MITRE ATLAS, or an OWASP Agentic
category), so the chain resolves to a named technique rather than a per-minister
default in every case; a per-minister default remains available as a fallback so
the chain is always populated, and the chain records which basis was used. The provenance record is serialized into the
SHA-256 hash-chained Ledger as a hashed field. Consequently, editing the
\emph{grounding} of a historical decision --- for instance, relabeling a
credential-access block as a benign read to mislead an auditor --- breaks the
hash chain at that entry, detected by \texttt{/ledger/verify}. Prior runtime
monitors provide either explainability or tamper-evident logging but not a
taxonomy-grounded chain whose integrity is cryptographically verifiable; the
ordered stage axis additionally provides the substrate for stage-aware
trajectory reasoning, which we leave to future work.

\subsection{OpenClaw Plugin}
\label{sec:plugin}

The Kavach plugin (\texttt{openclaw-plugin-kavach.ts}) implements the
OpenClaw plugin interface:

\begin{itemize}
  \item \textbf{\texttt{pre\_tool\_call}}: synchronously calls
        \texttt{/hook/parliament} before the agent executes a tool.
        On \textsc{Block} or \textsc{Escalate}, the plugin throws
        \texttt{KavachBlockError} or \texttt{KavachEscalateError}, preventing
        tool execution.
  \item \textbf{\texttt{on\_tool\_error}}: increments the circuit-breaker
        failure counter. After \texttt{kavachCircuitBreakerThreshold}
        consecutive Kavach server errors the plugin enters \emph{open-circuit}
        mode: all tool calls are blocked and an operator alert is raised.
\end{itemize}

The plugin is configured via \texttt{openclaw.plugin.json} with
\texttt{parliamentUrl} (default: \texttt{http://127.0.0.1:8088}) and
\texttt{SESSION\_ID} in the OpenClaw runtime environment.

We emphasize that Kavach itself is \emph{not} an OpenClaw component: all
detection logic --- embedding, the four Ministers, the Speaker, and the Ledger
--- lives in \texttt{parliament/server.py}, a standalone ASGI (FastAPI) service
exposing HTTP endpoints. The OpenClaw plugin is a $\sim$150-line adapter that
renders a tool call to text and \texttt{POST}s it; any agent framework that can
issue an HTTP request at a pre-execution hook integrates against the same
service without modification. The framework binding is an adapter, not a
dependency.

\subsection{Design Decisions and Trade-offs}

\paragraph{Ministers are domain-bound, not stage-bound.}
OpenClaw's five-stage pipeline (Init$\to$Execution) provides natural hooks but
creates stage-specific blind spots: an attack that appears in the Inference
stage looks different from the same attack in the Decision stage, requiring
per-stage classifiers. By binding Ministers to attack \emph{domains} we achieve
consistent coverage across all stages with a single collection per domain.

\paragraph{Append-only Ledger, no session objects.}
Storing session state in an append-only log rather than an in-memory object
prevents an attacker from clearing session state via a crafted tool call.
The immutability guarantee is enforced at the database level
(\texttt{kavach\_parliament.db} opened read-write only by the server; no
\texttt{DELETE} or \texttt{UPDATE} statements are issued).

\paragraph{Synchronous evaluation on the hot path.}
Kavach evaluates in the \texttt{pre\_tool\_call} hook synchronously.
Asynchronous evaluation would reduce latency but allow a race between the
verdict and execution; we accept the latency cost to maintain the security
guarantee that \emph{no tool executes before a verdict is issued}.
